\documentclass[UTF8, 11pt]{ctexart}

\usepackage[margin=1in]{geometry}
\usepackage{amsmath, amssymb}
\usepackage{booktabs}
\usepackage{hyperref}
\usepackage{enumitem}
\setlist{leftmargin=*}
\setlength{\parskip}{0.35em}
\setlength{\parindent}{0pt}

\title{\textbf{实验报告草稿}\\
分布式影响力最大化中的子模优化与贪心近似算法}
\author{Team Members: Name 1, Name 2, Name 3}
\date{CS240: 算法设计与分析，2026 年春季学期}

\begin{document}

\maketitle

\section{背景与相关工作}

影响力最大化是现代网络数据分析中的经典算法问题。给定一个社交网络、通信网络或多智能体交互网络，我们希望选择少量初始节点，使信息或行为经过网络传播后能够影响尽可能多的节点。该问题可以用于病毒式营销、公共信息传播、虚假信息干预，也可以被解释为多智能体网络中的有限预算决策问题。

从算法角度看，影响力最大化把现代网络应用转化为一个带基数约束的组合优化问题。目标函数通常是期望激活节点数。虽然直接求最优种子集合是 NP-hard 的，但在独立级联模型和线性阈值模型下，影响函数具有单调子模性，因此经典贪心算法具有 $1-1/e$ 的近似保证。

这一问题与分布式及大规模优化、多智能体决策与控制方向有自然联系。在大规模网络中，完整图结构往往无法由单个中心节点完全掌握；在多智能体系统中，每个 agent 只能访问局部邻居信息。如何在局部信息、通信约束和全局影响目标之间取得平衡，是分布式优化问题中的典型主题。

\section{代表性论文与方法分析}

本项目选择 Kempe, Kleinberg, and Tardos 关于影响力最大化的经典工作作为代表性论文。其形式化问题为：
\[
    \max_{S\subseteq V,\ |S|\leq k} \sigma(S),
\]
其中 $\sigma(S)$ 是种子集合 $S$ 在给定扩散模型下的期望影响范围。

核心算法是贪心选择。令当前种子集合为 $S_t$，第 $t+1$ 轮选择
\[
    v^\star
    =
    \arg\max_{v\in V\setminus S_t}
    \left(\sigma(S_t\cup\{v\})-\sigma(S_t)\right).
\]
然后令 $S_{t+1}=S_t\cup\{v^\star\}$，重复直到 $|S|=k$。

由于精确计算 $\sigma(S)$ 本身通常困难，实验中使用 Monte Carlo 模拟估计影响范围。若图有 $n$ 个节点、$m$ 条边、预算为 $k$、候选节点数约为 $n$、每次估计使用 $R$ 次模拟，则朴素贪心的运行代价约为
\[
    O(k n R (n+m)),
\]
其中一次独立级联模拟最多访问所有节点和边。

CELF 懒惰贪心算法利用子模函数的边际收益递减性质。每个候选节点的旧边际收益可以看作当前边际收益的上界，因此算法用优先队列维护候选节点，只有当队首节点的边际收益过期时才重新估计。这样通常能显著减少 Monte Carlo 估计次数。

\section{复现与实现计划}

项目将实现以下算法：
\begin{itemize}
    \item 随机选择基线；
    \item 度数中心性基线；
    \item 朴素贪心影响力最大化；
    \item CELF 懒惰贪心；
    \item 面向多智能体局部信息的分布式近似选择方案。
\end{itemize}

独立级联模型作为主要扩散模型。实验中，每条边 $(u,v)$ 具有传播概率 $p_{uv}$。当节点 $u$ 被激活后，它有一次机会以概率 $p_{uv}$ 激活邻居 $v$。传播过程停止后，激活节点总数即为一次模拟中的影响范围。

分布式扩展中，图被划分为若干 agent 持有的局部子图。每个 agent 计算本地候选节点的边际收益估计，并通过有限轮通信汇总候选节点。该扩展不要求达到中心化贪心的完整理论保证，而是用于展示分布式优化中的核心张力：局部计算降低集中式开销，但有限通信可能降低全局选择质量。

\section{实验评估计划}

实验将围绕以下指标展开：
\begin{itemize}
    \item \textbf{影响范围：} 最终平均激活节点数。
    \item \textbf{运行时间：} 不同算法在相同图和预算下的耗时。
    \item \textbf{模拟次数：} CELF 相对于朴素贪心减少了多少影响力估计。
    \item \textbf{边际收益曲线：} 随着预算 $k$ 增大，新增种子带来的增益是否递减。
    \item \textbf{分布式近似损失：} 与中心化贪心相比，局部信息方案的影响范围下降多少。
\end{itemize}

数据集包括 Erdos--Renyi 随机图、Barabasi--Albert 无标度图、Watts--Strogatz 小世界图和一个小型真实网络。实验将比较不同图规模、传播概率和预算设置下的算法表现。

\section{可选拓展}

若时间允许，项目将加入以下拓展：
\begin{itemize}
    \item 分析传播概率变化对最优种子结构的影响；
    \item 比较独立级联模型和线性阈值模型；
    \item 设计通信轮数受限的分布式贪心变体；
    \item 讨论该问题与分布式优化中的 consensus、局部目标聚合和近似收敛分析之间的联系。
\end{itemize}

\begin{thebibliography}{9}

\bibitem{kempe2003}
D. Kempe, J. Kleinberg, and E. Tardos.
\newblock Maximizing the spread of influence through a social network.
\newblock \emph{KDD}, 2003.

\bibitem{nemhauser1978}
G. L. Nemhauser, L. A. Wolsey, and M. L. Fisher.
\newblock An analysis of approximations for maximizing submodular set functions.
\newblock \emph{Mathematical Programming}, 1978.

\bibitem{leskovec2007}
J. Leskovec, A. Krause, C. Guestrin, C. Faloutsos, J. VanBriesen, and N. Glance.
\newblock Cost-effective outbreak detection in networks.
\newblock \emph{KDD}, 2007.

\end{thebibliography}

\end{document}
